\section*{Giới thiệu}
\addcontentsline{toc}{section}{Giới thiệu}
\textbf{Altium Designer} là phần mềm thiết kế điện tử tự động (EDA) dùng để thiết kế mạch in (PCB), FPGA, phần mềm nhúng và quản lý thư viện linh kiện. Mạch in (\textit{Printed Circuit Board}) vừa đỡ cơ khí, vừa nối điện các linh kiện bằng đường đồng (\textit{track}), \textit{pad} \dots\ được ăn mòn từ lớp đồng phủ trên đế cách điện.

\given{fig4_01}{0.7\linewidth}{PCB circuit in Altium Designer}

Trong bài này, thiết kế hai mạch ổn áp trên Altium Designer:
\begin{itemize}
  \item Mạch ổn áp tuyến tính dùng IC \textbf{7805} (mục~\ref{sec:7805});
  \item Mạch ổn áp xung hạ áp (\textit{buck converter}) dùng IC \textbf{LM2596} (mục~\ref{sec:lm2596}).
\end{itemize}
Mỗi mạch gồm hai phần: vẽ sơ đồ nguyên lý (\textit{Schematic}) và bố trí mạch in (\textit{PCB layout}: lớp TOP, lớp BOTTOM và ảnh 3D).

\begin{notebox}{Phím tắt thường dùng trong Schematic}
\textbf{Space}: xoay linh kiện \quad \textbf{X}: lật ngang \quad \textbf{Y}: lật dọc \quad \textbf{Ctrl + W}: đi dây (\textit{place wire})
\end{notebox}
